\clearpage
\section*{Memorandum to the Chief Operations Officer}
\addcontentsline{toc}{section}{Memorandum to the COO}
\noindent\textbf{To:} Chief Operations Officer, HiMCM\\
\textbf{From:} Cooling-service modeling team\\
\textbf{Subject:} A service-targeted cooling plan for the Fremont canopy corridor

For the modeled $8382$-m$^2$ occupied corridor beneath the Fremont Street Experience canopy \cite{fremont-canopy}, use \textbf{266 type-A fans and 16 type-C evaporative units} as the nominal planning configuration. At $38^\circ$C and $20\%$ RH, this is the minimum annual-cost integer plan meeting our global coverage, mean-relief, clearance and electricity limits. The four-hour attendance profile represents 3210 visitor-hours, with the largest crowd in the middle demand zone.

One optimal assignment puts 116 fans in zone 1, 148 fans and all 16 evaporative units in zone 2, and two fans in zone 3. It provides \textbf{81.79\% person-time coverage} and $1.5002^\circ$C mean reduction in SET excess. Zone coverages are $96.7\%$, $100\%$ and $1.7\%$: the global objective strongly favors the busiest sector. If service is required in every sector, add zonal coverage constraints before procurement. Equipment clearance occupies $1.87\%$ of the total service area; circulation routes and exits still require a layout check.

A four-hour operating period uses \textbf{426 kWh and approximately 2540 liters of make-up water}, delivering 4816 person$\cdot{}^\circ$C$\cdot$h of relief. Efficiency is $11.30$ such service units per kWh. Mean residual SET excess remains $6.27^\circ$C, so this configuration reduces heat burden without making the corridor thermally neutral.

The common price basis gives \textbf{USD 111,800 in equipment capital} and \textbf{USD 50,926/year} at 600 operating hours. The annual total comprises approximately USD 25,823 capital recovery, 11,180 maintenance, 12,780 electricity and 1,143 water. Obtain installation and utility quotes before purchase: the nominal counts persist in only 56 of 243 tested price combinations, and 69 combinations favor fans alone. A $1.4^\circ$C relief target instead selects 280 fans and nine C units at USD 45,502/year; a $1.6^\circ$C target selects 259 fans and 22 C units at USD 56,221/year.

Commission a representative section before full deployment. Measure air speed, temperature and humidity near, midway from and at the edge of each unit; confirm delivery to standing visitors, water supply and unobstructed movement. Update the allocation using the observed crowd distribution and conditions. The nominal mean-relief margin is very small, so the computed optimum offers no delivery reserve. Agree a margin or revised service target before implementation, and reoptimize when weather or attendance changes.

During severe heat, follow National Weather Service guidance: monitor forecasts, adjust strenuous outdoor activity, and provide breaks and access to cooler locations \cite{nws-heat}. Use official heat advisories to trigger measures beyond equipment operation.
